Se utilizó GCC 15.2.0 (MSYS2), C++17 con \texttt{-O2}, y Python 3.12 con NumPy, pandas y Matplotlib.
Se generaron 288 instancias con semilla 2026: filas $N=4$, $8$, $16$,
$64$, $256$, $1024$; columnas $M=N$ o $2N$; cuatro tipos, dos dominios
y tres muestras; $V=3$ y $P=10$ o $1000$. Un camino reservado sin
obstáculos garantiza que todas tengan solución.

El auxiliar \texttt{scripts/benchmark\_runner.cpp} mide cada algoritmo en
un proceso separado. Cronometra solo el solver, sus asignaciones y la
reconstrucción; usa la mediana de tres llamadas, salvo fuerza bruta en
$16\times16$, medida una vez. La memoria es el máximo
(\texttt{PeakWorkingSetSize}), con entrada, bibliotecas y
asignador~\cite{memoria_windows}, expresado en MiB.
Las figuras promedian las instancias de cada grupo
